% Part: lambda-calculus
% Chapter: syntax
% Section: alpha

\documentclass[../../../include/open-logic-section]{subfiles}

\begin{document}

\olfileid{lam}{syn}{alp}
\olsection{$\alpha$-转换}

$\lambd[x][x]$ 和 $\lambd[y][y]$ 之间有什么关系？它们都表示恒等函数。当然，它们是句法上不同的项。它们的区别仅在于约束变元的名称，而且其中一个可以视作对另一个进行约束变元“重命名”的结果。这被称为\emph{$\alpha$-转换}。

\begin{defn}[约束变元的更改，$\aconvone$]
  如果一个项 $M$ 中包含 $\lambd[x][N]$ 的一次出现，$y \notin \FV{N}$，且 $\Subst{N}{y}{x}$ 有定义，那么将此次出现替换为
  \begin{equation*}
    \lambd[y][\Subst{N}{y}{x}]
  \end{equation*}
  得到 $M'$，这一操作就称为\emph{约束变元的更改}，记作 $M \redone[\alpha] M'$。
\end{defn}

\begin{defn}[关系的兼容性]
  项上的一个关系 $R$ 被称为是\emph{兼容的}，如果它满足以下条件：
  \begin{enumerate}
  \item 如果 $R N N'$，那么 $R \lambd[x][N] \lambd[x][N']$
  \item 如果 $R P P'$，那么 $R (PQ) (P'Q)$
  \item 如果 $R Q Q'$，那么 $R (PQ) (PQ')$
  \end{enumerate}
\end{defn}

因此，让我们重新表述定义：

\begin{defn}[约束变元的更改，$\aconvone$]
  \emph{约束变元的更改}（$\redone[\alpha]$）是项上满足以下条件的最小兼容关系：
  \begin{align*}
    &\lambd[x][N] \redone[\alpha] \lambd[y][\Subst{N}{y}{x}] && \text{若
      $x \neq y$，$y \notin \FV{N}$} \\
    & &&\text{且 $\Subst{N}{y}{x}$ 有定义}
  \end{align*}
\end{defn}

这里的“最小”是指，该关系只包含由兼容性以及上述附加条件所要求的配对，而不包含任何其他配对。因此，这个关系也可以如下定义：

\begin{defn}[约束变元的更改，$\aconvone$] \ollabel{defn:aconvone}
  \emph{约束变元的更改}（$\aconvone$）归纳定义如下：
  \begin{enumerate}
  \item 如果 $N \aconvone N'$，那么 $\lambd[x][N] \aconvone
    \lambd[x][N']$ \ollabel{defn:aconvone1}
  \item 如果 $P \aconvone P'$，那么 $(PQ) \aconvone (P'Q)$ \ollabel{defn:aconvone2}
  \item 如果 $Q \aconvone Q'$，那么 $(PQ) \aconvone (PQ')$ \ollabel{defn:aconvone3}
  \item 如果 $x \neq y$，$y \notin \FV{N}$ 且 $\Subst{N}{y}{x}$ 有定义，那么
    $\lambd[x][N] \redone[\alpha] \lambd[y][\Subst{N}{y}{x}]$。
    \ollabel{defn:aconvone4}
  \end{enumerate}
\end{defn}

这些定义是等价的，但我们把证明留作练习。从现在开始，我们将使用归纳定义。

\begin{defn}[$\alpha$-转换，$\aconv$]
  \emph{$\alpha$-转换}（$\aconv$）是项上包含 $\aconvone$ 的最小自反且传递的关系。
\end{defn}

如上所述，“最小”意味着该关系只包含由传递性和 $\aconvone$ 所要求的配对，由此得到以下等价定义：

\begin{defn}[$\alpha$-转换，$\aconv$] \ollabel{defn:aconv}
  \emph{$\alpha$-转换}（$\aconv$）归纳定义如下：
  \begin{enumerate}
  \item 如果 $P \aconv Q$ 且 $Q \aconv R$，那么 $P \aconv R$。
    \ollabel{defn:aconv1}
  \item 如果 $P \aconvone Q$，那么 $P \aconv Q$。 \ollabel{defn:aconv2}
  \item $P \aconv P$。 \ollabel{defn:aconv3}
  \end{enumerate}
\end{defn}

\begin{ex}
  $\lambd[x][f x]$ 可以 $\alpha$-转换为 $\lambd[y][f y]$，反之亦然。
  非正式地说，它们都是接受一个参数并返回 $f$ 作用于该参数的函数，这里引用了环境变元~$f$。

  $\lambd[x][f x]$ 不能 $\alpha$-转换为 $\lambd[x][g x]$。
  非正式地说，它们分别引用环境变元 $f$ 和~$g$，这使它们成为不同的函数：在 $f$ 和 $g$ 取值不同的环境中，它们的行为不同。
\end{ex}

\begin{prob}
  下列各对项是 $\alpha$-可转换的吗？
  \begin{enumerate}
  \item $\lambd[x][\lambd[y][x]]$ 和 $\lambd[y][\lambd[x][y]]$
  \item $\lambd[x][\lambd[y][x]]$ 和 $\lambd[c][\lambd[b][a]]$
  \item $\lambd[x][\lambd[y][x]]$ 和 $\lambd[c][\lambd[b][a]]$
  \end{enumerate}
\end{prob}

\begin{lem}\ollabel{lem:fv-one}
  如果 $P \aconvone Q$，那么 $\FV{P} = \FV{Q}$。
\end{lem}

\begin{proof}
  对 $P \aconvone Q$ 的推导作归纳。
  \begin{enumerate}
  \item 如果最后应用的规则是 \olref{defn:aconvone4}，那么 $P$ 具有形式 $\lambd[x][N]$ 而 $Q$ 具有形式
    $\lambd[y][\Subst{N}{y}{x}]$，其中 $x \neq y$，$y \notin \FV{N}$ 且 $\Subst{N}{y}{x}$ 有定义。我们根据 $x \in \FV{N}$ 是否成立分情况讨论：
    \begin{enumerate}
    \item 若 $x \in FV(N)$，则：
      \begin{align*}
        \FV{\lambd[y][\Subst{N}{y}{x}]} &=
        \FV{\Subst{N}{y}{x}} \setminus \{y\} \\
        & = ((\FV{N} \setminus \{x\}) \cup \{y\}) \setminus \{y\}
         && \text{ 根据 \olref[sub]{thm:infv}} \\
        & = \FV{N} \setminus \{x\} \\
        & = \FV{\lambd[x][N]}
      \end{align*}
    \item 若 $x \notin FV(N)$，则：
      \begin{align*}
        \FV{\lambd[y][\Subst{N}{y}{x}]}
        & = FV{\Subst{N}{y}{x}} \setminus \{y\} \\
        & = \FV{N} \setminus \{x\}
         && \text{根据 \olref[sub]{thm:notinfv}} \\
        & = \FV{\lambd[x][N]}.
      \end{align*}
    \end{enumerate}
  \item 其他三种情况留作练习。
  \end{enumerate}
\end{proof}

\begin{prob}
  完成 \olref[lam][syn][alp]{lem:fv-one} 的证明。
\end{prob}

\begin{lem}\ollabel{lem:inv}
  如果 $P \aconvone Q$，那么 $Q \aconvone P$。
\end{lem}

\begin{proof}
  对 $P \aconvone Q$ 的推导作归纳。
  \begin{enumerate}
  \item 如果最后一条规则是 \olref{defn:aconvone4}，那么 $P$ 具有形式 $\lambd[x][N]$ 而 $Q$ 具有形式 $\lambd[y][\Subst{N}{y}{x}]$，其中 $x \neq y$，$y \notin \FV{N}$ 且 $\Subst{N}{y}{x}$ 有定义。首先，由 \olref[sub]{thm:clr} 我们有 $y \notin
    \FV{\Subst{N}{y}{x}}$。由 \olref[sub]{thm:inv} 可知，$\Subst{\Subst{N}{y}{x}}{x}{y}$ 不仅有定义，而且等于~$N$。那么由 \olref{defn:aconvone4}，我们有 $\lambd[y][\Subst{N}{y}{x}] \aconvone \lambd[x][\Subst{\Subst{N}{y}{x}}{x}{y}] = \lambd[x][N]$。
  \end{enumerate}
\end{proof}

\begin{prob}
  补全 \olref[lam][syn][alp]{lem:inv} 的证明。
\end{prob}

\begin{thm}
  α-转换是项上的一个等价关系，即它是自反的、对称的和传递的。
\end{thm}

\begin{proof}
  \begin{enumerate}
  \item 对每个项 $M$，通过\emph{零次}约束变元改名，$M$ 可变为 $M$。
  \item 如果 $P$ 通过一系列约束变元改名 α-转换为 $Q$，那么从 $Q$ 出发，只需（根据 \olref{lem:inv}）以相反顺序逆转这些改名即可得到~$P$。
  \item 如果 $P$ 通过一系列约束变元改名 α-转换为 $Q$，并且 $Q$ 通过另一系列改名转换为 $R$，那么我们可以先应用第一个系列再应用第二个系列，从而将 $P$ 变为 $R$。
  \end{enumerate}
\end{proof}

从现在起，我们说 $M$ 和 $N$ 是\emph{α-等价的}，记作 $M \aeq N$，当且仅当 $M$ α-转换为 $N$（正如我们刚刚证明的，这当且仅当 $N$ α-转换为~$M$）。

\begin{thm}\ollabel{thm:fv}
  如果 $M \aeq N$，那么 $\FV{M} = \FV{N}$。
\end{thm}

\begin{proof}
  由 \olref{lem:fv-one} 立刻可得。
\end{proof}

\begin{lem}\ollabel{lem:sub:R}
  如果 $R \aeq R'$ 且 $\Subst{M}{R}{y}$ 有定义，那么 $\Subst{M}{R'}{y}$ 有定义并且 α-等价于 $\Subst{M}{R}{y}$。
\end{lem}

\begin{proof}
  练习。
\end{proof}

\begin{prob}
  证明 \olref[lam][syn][alp]{lem:sub:R}。
\end{prob}

回顾在 \olref[sub]{sec} 中，代入在某些情况下是未定义的；然而，通过对项进行 α-转换，我们可以通过重命名约束变元来使得代入始终有定义。如下定理所示，结果保持 α-等价：

\begin{thm}\ollabel{thm:sub}
  对于任意 $M$、$R$ 和~$y$，存在 $M'$ 使得 $M \aeq M'$ 且 $\Subst{M'}{R}{y}$ 有定义。此外，如果存在另一对 $M'' \aeq M$ 和 $R''$，其中 $\Subst{M''}{R''}{y}$ 有定义且 $R'' \aeq R$，那么 $\Subst{M'}{R}{y} \aeq \Subst{M''}{R''}{y}$。
\end{thm}

\begin{proof}
  对 $M$ 的结构作归纳：
  \begin{enumerate}
  \item $M$ 是变元~$z$：练习。
  \item 假设 $M$ 具有形式 $\lambd[x][N]$。选取一个变元 $z$，它不同于 $x$ 和 $y$，并且满足 $z \notin \FV{N}$ 和 $z \notin \FV{R}$。由归纳假设，存在 $N'$ 使得 $N' \aeq N$ 且 $\Subst{N'}{z}{x}$ 有定义。那么根据 \olref{defn:aconvone}\olref{defn:aconvone1}，$\lambd[x][N] \aeq \lambd[x][N']$ 也成立。现在由 \olref{defn:aconvone}\olref{defn:aconvone4}，$\lambd[x][N']
    \aeq \lambd[z][\Subst{N'}{z}{x}]$。这是因为 $z \ne x$，$z \notin FV(N')$ 且 $\Subst{N'}{z}{x}$ 有定义。最后，因为 $z \neq y$ 且 $z \notin FV(R)$，所以 $\Subst{\lambd[z][\Subst{N'}{z}{x}]}{R}{y}$ 有定义。

    此外，如果存在另一对满足相同条件的 $N''$ 和 $R''$，
    \begin{multline*}
      \Subst{(\lambd[z][\Subst{N''}{z}{x}])}{R''}{y} =\\
    \begin{aligned}
      &= \lambd[z][\Subst{\Subst{N''}{z}{x}}{R''}{y}] \\
      &= \lambd[z][\Subst{\Subst{N''}{z}{x}}{R}{y}] && 由 \olref{lem:sub:R}\\
      &=\lambd[z][\Subst{\Subst{N'}{z}{x}}{R}{y}]
      && 由归纳假设\\
      &=\Subst{(\lambd[z][\Subst{N'}{z}{x}])}{R}{y}
    \end{aligned}
    \end{multline*}
  \item $M$ 具有形式 $(PQ)$：练习。
  \end{enumerate}
\end{proof}

\begin{prob}
  补全 \olref[lam][syn][alp]{thm:sub} 的证明。
\end{prob}

\begin{cor}\ollabel{cor:sub}
  对于任意 $M$、$R$ 和~$y$，存在 $M'$ 和 $R'$ 使得 $M' \aeq M$，$R \aeq R'$ 且 $\Subst{M'}{R'}{y}$ 有定义。此外，如果存在另一对 $M'' \aeq M$ 和 $R''$ 使得 $\Subst{M'}{R'}{y}$ 有定义，那么 $\Subst{M'}{R'}{y} \aeq
  \Subst{M''}{R''}{y}$。
\end{cor}

\begin{proof}
  由 \olref{thm:sub} 立刻可得。
\end{proof}

\end{document}